%% kapa: exam=1 round=29 year=2018 period=1 form=A subject=economics subject_ko=경제학원론 seq=27 number=67
%% source: Q-Net 감정평가사 기출문제 (cst00309) artlSeq=5202808 file=extracted/1008/5202808/2018년 제29회 감정평가사 1차시험 기출문제/2018년 제29회 감정평가사 1차시험 기출문제/1교시 A형.pdf
%% conversion: kordoc 4.24.51 -> tools/gen1.py (auto); review=proofread-sonnet; flags=none
\begin{question}{67}
경제정책에 관한 설명으로 옳은 것을 모두 고른 것은?
\begin{framedtext}
ㄱ. 외부시차는 경제에 충격이 발생한 시점과 이에 대한 정책 시행 시점 사이의 기간이다. \par
ㄴ. 자동안정화장치는 내부시차를 줄여준다. \par
ㄷ. 루카스(R. Lucas)는 정책이 변하면 경제주체의 기대도 바뀌게 되는 것을 고려해야 한다고 주장하였다. \par
ㄹ. 시간적 불일치성 문제가 있는 경우 자유재량적 정책이 바람직하다. \par
\end{framedtext}
\begin{choices}
  \item ㄱ, ㄴ
  \item ㄱ, ㄷ
  \item ㄱ, ㄹ
  \item ㄴ, ㄷ
  \item ㄴ, ㄹ
\end{choices}
\end{question}
\answer{4}
